\documentclass[10pt,a4paper]{amsart}
\usepackage[margin=19mm]{geometry}
\usepackage{amsmath,amssymb,mathtools}
\usepackage[scr=rsfs,cal=euler]{mathalfa}
\usepackage{xfrac}
\usepackage[unicode,hidelinks,bookmarks=false]{hyperref}
\newcommand{\ie}{i.e.\ }
\newcommand{\cf}{cf.\ }
\newcommand{\resp}{respectively\ }
\newcommand{\Subsection}{\subsection}
\providecommand{\nobreakdash}{\nobreak\hbox{-}}
\setlength{\emergencystretch}{2em}
\multlinegap0pt
\usepackage{bm}
\usepackage{bbm}
\usepackage{dsfont}
\usepackage{xspace}
\usepackage{cancel}
\usepackage{mathtools}
\usepackage{cleveref}
\usepackage{enumitem}
\usepackage{amsthm}
\makeatletter
\@ifundefined{theorem}{\newtheorem{theorem}{Theorem}}{}
\makeatother
\newcommand{\G}{\Gamma}
\newcommand{\calF}{\mathcal{F}}
\newcommand{\calG}{\mathcal{G}}
\newcommand{\calM}{\mathcal{M}}
\newcommand{\eps}{\epsilon}
\newcommand{\ov}[1]{\overline{#1}}
\newcommand{\ph}[1]{\phantom{#1}}
\newcommand{\rvline}{\hspace*{-\arraycolsep}\vline\hspace*{-\arraycolsep}}
\title[Nonunique Ergodicity on the Boundary of Outer space]{Nonunique Ergodicity on the Boundary of Outer space}
\author{Mladen Bestvina, Elizabeth Field, Sanghoon Kwak}
\date{Source: arXiv:2409.15591v2 (CC BY 4.0)}
\begin{document}
\maketitle

\noindent\emph{Source and licence.} Nonunique Ergodicity on the Boundary of Outer space. Version arXiv:2409.15591v2 (CC BY 4.0), distributed under the Creative Commons Attribution 4.0 International licence, as recorded in the arXiv licence metadata for this exact version and shown on the abstract page https://arxiv.org/abs/2409.15591v2. Full rights notice: licenses/arxiv-2409.15591-CC-BY-4.0.txt.\par\smallskip

\noindent\textit{Editorial scope.} Counted unit: the Theorem whose source label is \texttt{thm:lowerboundFinale\_unfolding} in the article (source0.tex lines 2288--2292, source label \texttt{thm:lowerboundFinale\_unfolding}), delivered together with the article's own complete original proof of it (source0.tex lines 2294--2403, one \texttt{proof} environment of 110 lines). No mathematics is written, repaired or re-derived: statement and proof are the article's own text line for line. Required context retained uncounted: thm:lowerboundFinale\_folding (lines 1621--1624); thm:TheGame (lines 2130--2132); thm:klp\_then\_ergodic (lines 2225--2237). Every definition, display and label of those lines is retained unchanged. Omitted from this package and not redistributed: 1--1620, 1625--2129, 2133--2224, 2238--2287, 2293--2293, 2404--3031. Nothing inside a retained excerpt is omitted or altered: every retained body is delivered exactly as the article's lines are written, so no omission receipt of any kind accompanies this package. Citations: the retained excerpts carry no citation command at all, so no bibliography entry of the article is transcribed and no reference is redistributed. Reference adaptation: none was needed and none was made; every reference inside the delivered unit resolves inside it, so no unresolved reference or citation is produced. Printed slips kept literal: no printed word, symbol, hypothesis or number of the counted statement or of its proof was altered. Layout and class: the article's own class is \texttt{amsart} with packages including geometry, alltt, float, marginnote, amsfonts, euscript, enumitem, latexsym, epsfig, verbatim, amsmath, amsthm, amssymb, colonequals; the wrapper is the lane's pinned \texttt{amsart} editorial wrapper at 10pt on A4, which restates the theorem-like environments the retained text needs and adds the article's own macro definitions that the retained text needs (\texttt{\textbackslash G}, \texttt{\textbackslash calF}, \texttt{\textbackslash calG}, \texttt{\textbackslash calM}, \texttt{\textbackslash eps}, \texttt{\textbackslash ov}, \texttt{\textbackslash ph}, \texttt{\textbackslash rvline}), with extra wrapper packages (bm, bbm, dsfont, xspace, cancel, mathtools, cleveref, enumitem). The delivered numbering is the unit's own. Assets: no retained span contains a figure environment, a graphics-inclusion command, a PDF-page inclusion, a table or a TikZ picture; no image file is copied into this package, the package declares no asset dependency and no image file is redistributed with it. Licence: version arXiv:2409.15591v2 of this article is distributed under the Creative Commons Attribution 4.0 International licence, as recorded in the arXiv licence metadata for this exact version and shown on the abstract page https://arxiv.org/abs/2409.15591v2. A full rights notice is delivered at licenses/arxiv-2409.15591-CC-BY-4.0.txt. Characters: every retained excerpt is pure ASCII, so the delivered engine, XeLaTeX with its default Latin Modern text font, reports no missing character.
\begin{theorem} \label{thm:lowerboundFinale_folding}
For each $n \ge 4$, there exists a reduced folding sequence $\calG=(G_k)_{k \ge 0}$ of
graphs of rank $n$ whose limiting tree admits at least $2n-6$ linearly independent ergodic length measures.
\end{theorem}

\begin{theorem}\label{thm:TheGame}
  Alice has a winning strategy. 
\end{theorem}

\begin{theorem} \label{thm:klp_then_ergodic}
Let $(G_s)_{s \le 0}$ be a reduced unfolding sequence with legal lamination $\Lambda \subset G_0$, and $\calM = (M^{r,s})_{r \le s \le 0}$ be the inverse system of $n \times n$ transition matrices associated with the unfolding sequence $(G_s)_{s \le 0}$. Let $1
\le m \le n$. Suppose that for each $s \le 0$, there exists a positive triple $(\ov{p}_s, \ov{\eps}_s,
\ov{\delta}_s)$ such that for every $r < s$, $M^{r,s} \in \calM$ is an $(m, \ov{p}_s,
\ov{\eps}_s,\ov{\delta}_s)$ matrix such that 
\begin{enumerate}[label=(\roman*)]
\item $\{\ov{p}_s\}_{s \le 0}$ is uniformly bounded,
\item $\lim_{s \to -\infty} \ov{\eps}_s =0$, and
\item $\lim_{s \to -\infty} \ov{\delta}_s =0$.
\end{enumerate}
Then the legal
lamination $\Lambda$ admits at least $m$ linearly independent ergodic measures.
\end{theorem}

\begin{theorem} \label{thm:lowerboundFinale_unfolding} 
For each $n \ge 4$, there exists a reduced unfolding sequence
$\calG=(G_k)_{k \le 0}$ of graphs
of rank $n$ whose legal lamination admits at least $2n-6$ linearly independent ergodic measures.
\end{theorem}

\begin{proof}

Following the construction in \Cref{thm:lowerboundFinale_folding}, build the
same graph $\G_{a_0}$ with the maps $\calF^{(k)}$, to form a reduced unfolding sequence
$\calG = (G_k)_{k\le 0}=\{\phi_{k-1} : G_{k-1} \to G_{k}\}_{k \le 0}$ of graphs
of rank $n$ with $G_k = \G_{a_0}$ and $\phi_k = \calF^{(k)}$. The only
difference we will make for our current setting of an unfolding sequence is that
we will label the edges of $\Gamma_{a_0}$ so that
\[
(e^1,\ldots,e^{2n-3})=(a_{n-3},a_{n-2},\ldots, a_1, b_{n-3},b_{n-2},\ldots, b_1,a_0,b_0,c).
\] 
Then, our
transition matrix $M$ for $\calF$ will be of the form:
\begin{center}
  \scalebox{0.9}{$
    \begin{pmatrix}
    \setlength\arraycolsep{5pt}
    \begin{matrix}
        2\alpha_{n-3}  &  3             &  2                 & 2          & \cdots    & 2 \\
                       &  2\alpha_{n-4} &  3                 & 2          & \cdots    & 2 \\
                       &                &  \ddots            & \ddots     & \ddots    & \vdots \\
                       &                &                    & 2\alpha_3  & 3         & 2 \\
                       &                &  \ph{\alpha_{n-5}}      &            & 2\alpha_2 & 3 \\
                       &                &                    &            &           & 2\alpha_1
    \end{matrix}
    &
    \rvline %vertical separater
    &
    \setlength\arraycolsep{8pt}
    \begin{matrix}
        2                         &  2                 &   2       &   2        &   \cdots  & 2 \\
        \ph{\beta_{n-3}}          &  2                 &   2       &   2       &   \cdots  & 2 \\
                                  & \ph{\beta_{n-4}}   &   \ddots  &   \ddots   &   \ddots  & \vdots \\
                                  &                    &           &   2        &    2      & 2 \\
                                  &                    &           &            &    2      & 2 \\
                                  &                    &           &            &           & 2
    \end{matrix}
    &
    \rvline
    &
    \begin{matrix}
      0 & 0      & 2 \\
      0 & 0      & 2 \\
  \vdots & \vdots & \vdots \\
      0 & 0      & 2 \\
      0 & 0      & 2 \\
      1 & 0      & 2
    \end{matrix}\\
    \hline %%%%%%%END OF FIRST BLOCK ROW
    \setlength\arraycolsep{7pt}
    \begin{matrix}
          \alpha_{n-3} &   \alpha_{n-4}       &   3             & 3              & \cdots         &     3 \\
                       &   \alpha_{n-4}       &   \alpha_{n-5}  & 3              & \cdots         &     3 \\
                       &                      &   \ddots        & \ddots         & \ddots         &     \vdots \\
                       &                      &                 & \alpha_3       & \alpha_2       &     3 \\
                       &                      &                 &                & \alpha_2       &     \alpha_1 \\
                       &                      &                 &                &                &     \alpha_1
    \end{matrix}
    &
    \rvline %vertical separater
    &
    \setlength\arraycolsep{6.5pt}
    \begin{matrix}
       \ph{2}\beta_{n-3}    &  4             &   3                    & 3          & \cdots         & 3 \\
                       &  \beta_{n-4}   &   4                    & 3          & \cdots         & 3 \\
                       &                &   \ddots               & \ddots     & \ddots         & \vdots \\
                       &                & \ph{\beta_{n-3}}       & \beta_3    & 4              & 3 \\
                       &                &                        &            & \beta_2        & 4 \\
                       &                &                        &            &                & \beta_1
    \end{matrix}
    &
    \rvline
    &
    \begin{matrix}
      0 & 0      & 3 \\
      0 & 0      & 3 \\
  \vdots & \vdots & \vdots \\
      0 & 0      & 3 \\
      0 & 0      & 3 \\
      0 & 1      & 0
    \end{matrix}\\
    \hline
    \setlength\arraycolsep{9pt}
    \begin{matrix}
        1            & \ph{\alpha}0\ph{\alpha_n} & 0 & \ph{\alpha}0 & \cdots & 0\\
        \alpha_{n-3} & \ph{\alpha}3\ph{\alpha_n} & 3 & \ph{\alpha}3 & \cdots & 3\\
        1            & \ph{\alpha}1\ph{\alpha_n} & 1 & \ph{\alpha}1 & \cdots & 1
    \end{matrix}
    &
    \rvline
    &
    \setlength\arraycolsep{8.3pt}
    \begin{matrix}
       \ph{\beta}0\ph{\beta} & \ph{\beta_n}0\ph{\beta} & \ph{1}0 & \ph{\beta}0 & \cdots & 0\\
       \ph{\beta}4\ph{\beta} & \ph{\beta_n}3\ph{\beta} & \ph{1}3 & \ph{\beta}3 & \cdots & 3\\
       \ph{\beta}1\ph{\beta} & \ph{\beta_n}1\ph{\beta} & \ph{1}1 & \ph{\beta}1 & \cdots & 1
    \end{matrix}
    &
    \rvline
    &
    \begin{matrix}
        0 & 0 & 0\\
        0 & 0 & 3\\
        0 & 0 & 1
    \end{matrix}
    \end{pmatrix}
    $}
  \end{center}
 which can be seen as Bob's $(2n-6, p=2, \eps_r, \delta_r)$-matrix with appropriate choice of $(\alpha_1,\ldots,\allowbreak\alpha_{n-3},\beta_1,\ldots,\beta_{n-3})$ for each of Alice's choices of $\eps_r >0$ and $\delta_r>0$. Hence, Alice's winning sequence $\{(\eps_r, \delta_r)\}_{r \le 0}$ in \Cref{thm:TheGame} satisfying the conditions (i--iii) of \Cref{thm:klp_then_ergodic} provides at least $m=2n-6$ linearly independent ergodic measures on the legal lamination of our reduced unfolding sequence $\calG$.
\end{proof}

\end{document}
